% Counting Operator
\section{Counting Operator}

\subsection{Definition}

The counting operator $\text{COUNT}(n)$ assigns a cardinal value $n \in \mathbb{N}$ to an entity. This operator implements finite arithmetic on discrete quantities.

\subsection{Formal Characteristics}

\begin{itemize}
    \item \textbf{Position}: Medial or post-instantiation
    \item \textbf{Cardinality}: May repeat for multiple quantities
    \item \textbf{Dependency}: Requires prior instantiation (entity to count)
    \item \textbf{Notation}: Unary stroke system (tally marks)
\end{itemize}

\subsection{Stroke System}

\subsubsection{Basic Counting Signs}
The vertical stroke signs (M328--M335) form a unary numeral system:
\begin{align*}
    | &\rightarrow 1 \\
    || &\rightarrow 2 \\
    ||| &\rightarrow 3 \\
    |||| &\rightarrow 4 \\
    \vdots
\end{align*}

\subsubsection{Maximum Observed Values}
\begin{itemize}
    \item Most common: 1--7 strokes
    \item Rare: 8--12 strokes
    \item Practical limit: approximately 12 strokes
\end{itemize}

\subsection{Operational Semantics}

The counting operator expresses:
\begin{enumerate}
    \item \textbf{Cardinality}: How many of the instantiated entity
    \item \textbf{Quantity verification}: Administrative control value
    \item \textbf{Transaction amount}: Commercial quantity
\end{enumerate}

\subsection{Positional Patterns}

Counting signs follow strict positioning:
\begin{itemize}
    \item [INST][COUNT] = Standard pattern (entity + quantity)
    \item [INST][REL][COUNT] = Complex pattern (entity + relation + quantity)
    \item *[COUNT][INST] = Unattested (counting precedes instantiation)
\end{itemize}

The prohibition on count-before-instantiation is a strong constraint supporting the operational interpretation.

\subsection{Mathematical Properties}

The counting system exhibits:
\begin{itemize}
    \item \textbf{Closure}: Limited range (1--12 typically)
    \item \textbf{No zero}: Absence of entity rather than zero quantity
    \item \textbf{No negative values}: Cardinal-only system
    \item \textbf{No fractions}: Discrete units only
\end{itemize}

This is characteristic of administrative record-keeping rather than general arithmetic.

\subsection{Comparative Systems}

Similar counting systems appear in:
\begin{itemize}
    \item Egyptian unary numerals
    \item Roman numerals (I, II, III, IV...)
    \item Chinese rod numerals
    \item Tally sticks in medieval Europe
\end{itemize}
